\section{Encoder Implementation and Performance}
\label{sec:implementation}
\label{sec:transposed-comparison}

We implement the certified BCH--IMT encoders from
Section~\ref{sec:finite-certificates}, with an additional small-length variant
described in Appendix~\ref{app:small-k-variant}. The transposed encoder maps $N$
input blocks to $K$ output blocks; each $128$-bit block carries $128$
parallel binary instances. It applies the transposed inner, reverses the
routing, and applies the transposed BCH outer.

The configurations share an IMT template and fixed XOR circuits for the
inner maps. Small instances use direct routing; larger instances use
tiled routing. Packed indices and workspace-local page routing reduce
data movement. Online calls allocate no memory.
Symbolic checks, dense reference implementations, and adjoint tests check
the optimized maps. Appendix~\ref{app:implementation} gives the kernel
details, memory costs, correctness checks, and source-to-certificate bindings.

\paragraph{Transposed Comparison.}
Table~\ref{tab:transposed-comparison} compares binary encoders on one
Ryzen 9 7950X thread, pinned to CPU~15.
We compile with GCC~15.2.0, release optimization, and Zen~4 tuning,
with boost disabled; compiler flags differ per implementation
(Appendix~\ref{app:implementation}). The rate-$1/2$ BCH circuit uses AVX-512.
Except for the original-BAA reference, which is the authors' measurement
of the original BAA implementation, each cell is the median of three
serial process medians; each process has three warmups and $101$ timed calls.
Setup, allocation, and output checks are excluded. Input buffers are
initialized once, with no copy or reset between calls. These are
steady-state encoder latencies.

We include original BAA~\cite{KolesnikovEtAl2026BA} and the newer
chosen-block Golay and RM implementations~\cite{cryptoeprint:2026/1903}.
The chosen-block message lengths are $2^m-64$ for
$m\in\{16,18,20\}$; the other rows use $K=2^m$.
Expand--Convolute (EC) uses the conservative profile
$(w,m)=(33,25)$~\cite{RaghuramanRindalTanguy2023EC}, with its weaker
distance and margin. Binary RAA uses repetition three or four
\cite{blaze25}. Routing is included: BAA and RAA use Feistel permutations,
whereas SPIN uses packed tables. The SPIN cells use the certificate-matched
configurations under the same host and timing policy.
Rate-$1/2$ SPIN uses $(t,s)=(64,12)$ with two transvections per step at
$K=2^{16}$, and the original $(128,19)$ inner at the larger lengths.
All SPIN setup is precomputed; these timings do not use heuristic code refresh.

\begin{table}[tb]
\centering
\caption{Same-host transposed-encoder latency in milliseconds.
Entries are medians of three process medians, with no input resets,
except for the approximate original-BAA reference.
The pair $(\delta,\lambda)$ gives the relative distance and margin
in bits at the largest listed length, except for the RAA reference marked
$*$ below; -- denotes no bound or no timing reported.
EC uses the conservative, weaker distance and margin target.}
\label{tab:transposed-comparison}
\small
\setlength{\tabcolsep}{3pt}
% Generated by paper/build_imt_comparison.py; do not edit.
\begin{tabular}{@{}lccrrr@{}}
\toprule
Encoder & Rate & $(\delta,\lambda)$ & $K\approx2^{16}$ & $K\approx2^{18}$ & $K\approx2^{20}$ \\
\midrule
SPIN, IMT & $1/2$ & $(0.10,50.062)$ & 0.340 & 1.460 & 9.289 \\
Original BAA & $1/2$ & $(0.10,20)$ & -- & -- & ${\approx}32$ \\
Chosen Golay BAA & $1/2$ & $(0.10,46.386)$ & 0.595 & 2.781 & 23.997 \\
Chosen RM BAA & $1/2$ & $(0.10,47.093)$ & 0.866 & 3.807 & 27.660 \\
EC $(33,25)$ & $1/2$ & $(0.05,20)$ & 3.189 & 13.128 & 181.825 \\
Binary RAA, repetition 3 & $1/3$ & $\text{--}$ & 2.834 & 14.171 & 87.643 \\
Binary RAA, repetition 4 & $1/4$ & $(0.19,{\approx}13)^{*}$ & 1.009 & 6.643 & 45.176 \\
\bottomrule
\end{tabular}

\par\smallskip
\begin{minipage}{\linewidth}
\footnotesize
$*$RAA reference bound at $K=2^{22}$, not the timed $K=2^{20}$:
approximately $13$ bits without setup testing~\cite[Sec.~3.1]{blaze25}.
One-time testing of low-weight messages can improve the margin to about
$40$ bits ($41.5$ bits at this size), with near-quadratic setup work
\cite[Fig.~2]{blaze25}. No such test is used in our benchmark.
\end{minipage}
\end{table}

At $K=2^{20}$, rate-$1/2$ SPIN takes $9.289$\,ms, versus $23.997$\,ms
and $27.660$\,ms for chosen Golay and RM BAA.
The original rate-$1/2$ BAA baseline takes approximately $32$\,ms,
giving the approximately $3.4\times$ comparison in the abstract.
The table reports each construction's own distance guarantee, not a
common application-security level. In particular, the RAA reference
bound is not transferred to our smaller timed lengths.

\paragraph{Other Selected Configurations.}
Rate-$1/4$ SPIN takes $15.254$\,ms at $K=2^{20}$.
Rate-$1/2$ ordinary encoding takes $9.031$\,ms at that same $K$, under
the same precomputed, repeated-call timing policy. Thus the two directions
have similar cost.
Appendix~\ref{app:implementation} gives the full tables and protocols.
